% ============================================================================
%  frontmatter/resume.tex — Résumé FR + Abstract EN
%  À rédiger EN DERNIER (J12), une fois les résultats connus.
%  Cible : 200 à 250 mots chacun. Structure IMRaD condensée :
%  contexte -> problème -> méthode -> résultats chiffrés -> portée.
% ============================================================================
\chapter*{Résumé}
\addcontentsline{toc}{chapter}{Résumé}
\markboth{Résumé}{Résumé}

Le développement du firmware des produits industriels embarqués se heurte à des
cycles d'intégration longs, à des environnements de compilation hétérogènes et
à des exigences de cybersécurité croissantes. Ce mémoire présente la conception et la mise
en œuvre d'une infrastructure d'intégration
et de livraison continues pour le firmware du module de communication d'un
compteur d'énergie (\emph{Power Meter}) développé chez \Entreprise, module
chargé d'intégrer les protocoles Modbus, Ethernet, BACnet et une API REST.

Le travail ne part pas d'une feuille blanche : il reprend une chaîne GitHub
Actions existante, en analyse les limites, formalise des exigences avec
l'équipe, puis la fait évoluer. Les évolutions portent sur une image de
compilation unique intégrant le kit de développement de la cible, une matrice
de construction native et croisée, la mesure de couverture, les tests
d'intégration, l'analyse de composition binaire, la publication des artefacts
et la collecte d'indicateurs vers VictoriaMetrics et Grafana.

L'évaluation porte sur 164 exécutions sur cinq semaines. La chaîne applicative
rend un résultat complet en 21,8 minutes en médiane, la parallélisation de la
matrice réduisant sa durée d'environ un tiers. La chaîne restitue à chaque
exécution la couverture (32,3~\% des lignes), les défauts d'analyse statique et
l'inventaire des composants tiers avec leurs vulnérabilités. La distribution
Linux reste reconstruite sans cache en 68,2 minutes, ce qui fait du cache partagé
la principale perspective.

\vspace{1cm}
\noindent\textbf{Mots-clés :} intégration continue, livraison continue, DevOps,
systèmes embarqués, Yocto, analyse statique, chaîne d'approvisionnement
logicielle, analyse de composition, observabilité.

\vfill

\section*{Abstract}
\markboth{Abstract}{Abstract}

The development of firmware for embedded industrial products is hindered by
long integration cycles, heterogeneous build environments, and growing
cybersecurity requirements. This
thesis presents the design and implementation of a continuous integration and
continuous delivery
infrastructure for the firmware of the communication module of an energy
meter (\emph{Power Meter}) developed at \Entreprise, a module responsible for
integrating the Modbus, Ethernet, BACnet, and REST API protocols.

Rather than starting from scratch, the work takes over an existing GitHub
Actions pipeline, analyses its limits, formalises requirements with the team
and evolves it: a single build image embedding the target SDK, a native and
cross-compilation build matrix, code coverage, integration tests, binary
composition analysis, artifact publication, and metrics collection into
VictoriaMetrics and Grafana.

The evaluation covers 164 runs over five weeks. The application pipeline
delivers a complete result in a median of 21.8 minutes, the parallel matrix
cutting its duration by about one third. Each run reports coverage (32.3\% of
lines), static analysis defects and the inventory of third-party components
with their vulnerabilities. The Linux distribution is still rebuilt without
cache in 68.2 minutes, which makes a shared build cache the main next step.

\vspace{1cm}
\noindent\textbf{Keywords:} continuous integration, continuous delivery, DevOps,
embedded systems, Yocto, static analysis, software supply chain,
software composition analysis, observability.

\clearpage
